% !PRINTED 244
\weberpage{244}
%<<CONT
where $c_{h,k}$ are the coefficients to be determined, which by 3. satisfy the condition
\[
c_{h,k} = c_{k,h}\text{.}
\]
To determine them, we set in (29)
\[
x = u, \quad y = v,
\]
whereby $F_p$ vanishes, and obtain
\weq{30}{(u^p - v)(u - v^p) = \sum^{h,k} c_{h,k} u^h v^k.}
From this it follows first of all that
\[
c_{p,p} = 0
\]
must hold, since by (28), in the development (of the left side) in powers of $q$, the power $q^{-2(p^2+p)}$ cannot occur, and we can now put (29) into the form:
\weq{31}{(u^p - v)(u - v^p) = \sum_{0,p}^{h} \sum_{0,h-1}^{k} c_{h,k}(u^h v^k + u^k v^h) + \sum_{0,p-1}^{h} c_{h,h} u^h v^h.}
In this there are to be inserted the developments, resulting from (25) and (27), of
\[
u^h v^k + u^k v^h, \quad u^h v^h
\]
in powers of $q$, whose initial terms
\[
q^{-2(hp+k)}, \quad q^{-2h(p+1)}
\]
have the coefficients 1.
On the right side of (31) no two terms occur whose development begins with the same power of $q$; for from
\[
hp + k = h'p + k'
\]
it follows that $k \equiv k'\ (\mathrm{mod}\ p)$, and hence, since $k$ and $k' < p$, that $k = k'$, $h = h'$.
If one therefore arranges the series which represent the two sides of (31) according to ascending powers of $q$, and then sets equal to one another the coefficients of equally high powers, one obtains a sequence of linear equations for the unknowns $c_{h,k}$, of which each succeeding one contains only \gesperrt{one} new unknown, and this with the coefficient 1. The known terms of these equations which arise from the left side are, by (28), all \gesperrt{integers divisible by $p$}, and hence the $c_{h,k}$ likewise have integer values divisible by $p$. Accordingly we have
\weq{32}{F_p(x, y) = (x^p - y)(x - y^p) - p \sum_{0,p}^{h,k} a_{h,k} x^h y^k,}
%CONT>>
